\begin{table}[H]
\centering
\caption{Stage-4 Moment Blocks and Inputs}
\label{tab:struct_target_blocks}
\begin{threeparttable}
\footnotesize
\begin{tabular}[t]{>{\raggedright\arraybackslash}p{2.6cm}>{\raggedright\arraybackslash}p{3.5cm}>{\raggedright\arraybackslash}p{2.8cm}>{\raggedright\arraybackslash}p{4.7cm}}
\toprule
Block & Moments / inputs & Criterion weight & Role in estimation \\
\midrule
Observed treatment effects & 3 market ITTs & inverse variance (54.9--356.0) & Fit ATEs while residuals are shrinkage-regularized. \\
Social-channel diagnostics & 3 peer/rank composites & inverse diagnostic variance (100.0--200.0) & Discipline the common peer/rank scale. \\
Sorting first stage & 3 within-class dispersion changes & scale weight (1000.0) & Separate sorting compression from assignment payoff. \\
Revealed assignment payoff & Kenya treatment $\times$ predicted assignment-gain slope & inverse variance (216.3) & Pin down delivery activation at realized fidelity. \\
Fixed primitives and environment & $(\hat{\rho},\hat{\omega},\hat{\tau})$; peer shifts, class size, grade dispersion & not criterion targets & Measured inputs for observed and counterfactual regimes. \\
\bottomrule
\end{tabular}
\begin{tablenotes}[para,flushleft]
\footnotesize
\item \textit{Notes:} The stage-4 minimum-distance criterion targets only moments with a criterion weight in this table. Classroom peer shifts, class size, grade dispersion, and the first-three-block primitives enter the production mapping as measured inputs, not as additional fitted moments. This distinction prevents the target list from overstating what the outcome block estimates.
\end{tablenotes}
\end{threeparttable}
\end{table}
